\newpage
\subsubsection{Viabilidad económica del prototipo}

El análisis financiero del proyecto demuestra que el desarrollo del
prototipo es económicamente viable dentro de las condiciones y recursos
disponibles del equipo. La Tabla~\ref{tab:resumen_costos} presenta
el resumen consolidado de los costos totales estimados.

\begin{longtable}{|>{\raggedright\arraybackslash}m{8cm}|>
{\centering\arraybackslash}m{5cm}|}
\caption{Resumen de costos totales del prototipo.}
\label{tab:resumen_costos} \\
\hline
\textbf{Categoría} & \textbf{Total estimado (MXN)} \\
\hline
\endfirsthead
\multicolumn{2}{c}{{\bfseries Tabla \thetable\ -- continuación}} \\
\hline
\textbf{Categoría} & \textbf{Total estimado (MXN)} \\
\hline
\endhead

Infraestructura física (hardware) & \$8,000.00 \\
\hline
Transporte y logística & \$240.00 \\
\hline
Consumo eléctrico & \$2,720.00 \\
\hline
Licencias de software & \$0.00 \\
\hline
Servicios en la nube (AWS Free Tier) & \$0.00 \\
\hline
\textbf{Total general} & \textbf{\$10,960.00 MXN} \\
\hline
\end{longtable}

El costo total estimado del prototipo asciende a \$10,960.00 MXN,
distribuido entre tres integrantes del equipo, lo que representa
aproximadamente \$3,653.33 MXN por integrante. Este monto es accesible
para un proyecto académico de titulación, considerando que los costos
más significativos corresponden a adquisiciones de hardware que ya
han sido realizadas.

Los factores que contribuyen a la viabilidad económica del proyecto son:

\begin{itemize}
    \item \textbf{Cero costo en servicios en la nube:} La capa gratuita
    de AWS cubre el 100\% de los servicios de infraestructura requeridos
    durante la fase de prototipo, eliminando uno de los principales
    rubros de gasto en proyectos de esta naturaleza.
    \item \textbf{Cero costo en licencias de software:} Todas las
    herramientas de desarrollo seleccionadas son de código abierto o
    cuentan con versiones gratuitas para uso académico.
    \item \textbf{Dataset propio:} La recolección del conjunto de
    imágenes se realizará directamente en el Bosque de Chapultepec
    con el dron y los dispositivos del equipo, evitando la adquisición
    de datasets de pago.
    \item \textbf{Hardware reutilizable:} Tanto la Raspberry Pi como
    el dron son recursos que podrán seguir siendo utilizados más allá
    del prototipo, por lo que su costo representa una inversión con
    valor residual.
\end{itemize}

En caso de que el sistema evolucione hacia un ambiente productivo,
los costos de infraestructura aumentarían proporcionalmente al volumen
de usuarios y datos gestionados. El análisis de costos en ambiente
productivo y la proyección financiera a largo plazo se presentan en
la sección~\ref{subsec:produccion}.